\chapter{生态链：LP、GP、被投企业与服务商}

\section{LP 与 GP：信托关系与经济条款}

有限合伙人（Limited Partner, LP）是基金的主要出资方，以认缴承诺为上限承担有限责任；普通合伙人（General Partner, GP）或其关联实体担任基金管理人，负责投资决议、投后管理与投资者关系。二者关系由\textbf{有限合伙协议（LPA）}及附属 side letter 约束，核心经济条款包括管理费基数与费率、业绩报酬（carried interest）与门槛收益率（hurdle）、分配瀑布（European vs American waterfall）以及 GP 跟投（GP commit）比例。法律上 LP 通常不得干预具体项目投资，以免丧失有限责任保护；实务中通过咨询委员会（advisory committee）与协议约定的重大事项同意权行使制衡。

\section{基金载体、管理人与被投企业}

常见结构为：一支（或数支并列）\textbf{有限合伙基金}作为投资主体，向被投企业注资；\textbf{基金管理公司}（常由 GP 关联方设立）向基金收取管理费并聘用投资团队。被投企业处于生态链的「资产端」：其董事会构成、信息披露节奏与股东协议中的保护条款，直接决定 GP 能否有效行使治理。理解这一\textbf{三角结构}有助于划分「基金层面的义务」与「管理公司层面的义务」，例如在跨境架构中税务与监管申报的责任主体。

\section{交易服务商的角色}

\textbf{财务顾问（FA）}组织买方或卖方流程、协调尽调数据室与谈判节奏，但不替代 GP 的投资判断。\textbf{律师事务所}起草交易文件与 LPA，并出具法律意见。\textbf{审计师}提供基金与被投企业的财务审计或商定程序报告，支撑估值与 LP 报告。\textbf{基金行政管理人（fund administrator）}处理 NAV、投资者报表与监管报送接口。\textbf{托管行}在适用法域下保管现金与证券。这些角色嵌入在「立项—签约—交割—季报—退出」的时间轴上；冷启动阶段可用外包与里程碑付款降低固定成本，但随着 AUM 与 LP 要求上升，服务商层级与独立性往往需加强。

\section{网络效应与声誉资本}

生态链不是静态供应链：同一批 LP、FA 与律师在多个基金周期中重复博弈，\textbf{声誉}与\textbf{执行记录}成为比单次条款更重要的资产。新 GP 在构建生态位时，应优先建立可验证的里程碑（透明报告、按时交割、清晰沟通），以便在后续轮次中获得更优的介绍链与条款空间。
